\section{Prefix agreement}
\label{sec.analysis}

Throughout this section the execution is $(\epsilon,\lambda)$-typical and
\Cref{asm.majority} holds. We use the constants
\begin{gather*}
    \advrone = \frac{\advr}{1-f} + \frac{f_2}{f} + 2\epsilon, \qquad
    \advrtwo = \frac{\advrone}{1-\epsilon}, \\
    \slack = 2\epsilon + \frac{\advr f}{1-f}, \qquad
    \maxr = 1 + \advr + \slack,
\end{gather*}
and we write $k_0 = 2(\lambda + 2)f$. Since $\lambda \geq 2 / f$ we have
$2\lambda f \leq k_0 \leq 4\lambda f$, so $k_0$ is of the same order as the
common prefix parameter of~\cite{gkl24}. As $f$ and $\epsilon$ tend to $0$,
$\advrone$ and $\advrtwo$ tend to $\advr$, since $f_2/f<f/(2(1-f)^2)$, and $\maxr$
tends to $1 + \advr = n / (n - t)$.

\begin{lemma}\label{lem.rho}
    $\advrtwo<1$.
\end{lemma}
\begin{proof}
    By definition, $\advrtwo<1$ is equivalent to
    $\frac{\advr}{1-f}+\frac{f_2}{f}<1-3\epsilon$. From \Cref{asm.majority},
    $\advr\le1-\delta$ and $f<1/3$, and $f_2/f<f/(2(1-f)^2)$. So it suffices to
    show
    \[
        \frac{1-\delta}{1-f}+\frac{f}{2(1-f)^2} < 1-3\epsilon .
    \]
    Multiplying by $(1-f)^2$, this is
    $1-\tfrac{f}{2}-\delta(1-f)<(1-3\epsilon)(1-f)^2$, that is,
    $\delta(1-f)>\tfrac{3f}{2}-f^2+3\epsilon(1-f)^2$. Since $f<1/3$, the right
    side is at most $(1-f)(2f+3\epsilon)$. So $\delta>2f+3\epsilon$ suffices,
    and it follows from $3f+3\epsilon<\delta$.
\end{proof}

\begin{lemma}\label{lem.honestone}
    In a typical execution, a chain contains at most one honest block computed
    in any given round.
\end{lemma}
\begin{proof}
    Suppose towards a contradiction that $\ch[j]$ and $\ch[j']$ with $j<j'$ are
    both honest and both computed at round $i$. The honest validator that
    computed $\ch[j']$ extended at round $i$ the chain it adopted at round $i$.
    A block determines all of its ancestors, so that chain is $\ch[{:}j']$, and
    it contains $\ch[j]$. From \Cref{lem.membership}, $\ch[{:}j']\in\bt_{i-1}$,
    so $\ch[j]$ was computed at a round at most $i-1$. This contradicts the
    assumption that $\ch[j]$ was computed at round $i$.
\end{proof}

\begin{lemma}\label{lem.length}
    In a typical execution the following hold for every integer $v\ge0$.
    \begin{enumerate}
        \item Every chain $\ch\in\bt_v$ satisfies
            $\len{\ch}-1<\maxr f\max(v,\lambda)$.
        \item If $v>\lambda$, every chain $\ch$ adopted by an honest validator
            at round $v$ satisfies $\len{\ch}-1>(1-\epsilon)f(v-1)$.
    \end{enumerate}
\end{lemma}
\begin{proof}
    Let $w=\max(v,\lambda)$ and let $\ch\in\bt_v$. The blocks of $\ch[1{:}]$
    were computed in rounds of $[1,w]$. From \Cref{lem.honestone}, the honest
    ones were computed in distinct rounds, each of them successful. So there
    are at most $X([1,w])$ honest blocks, and at most $Z([1,w])$ adversarial
    ones. Since $w\ge\lambda$, \Cref{lem.typical} applies and
    \[
        \len{\ch}-1 = \len{\ch[1{:}]} <
            (1+\epsilon)fw+\left(\frac{\advr}{1-f}+\epsilon\right)fw = \maxr f w,
    \]
    where the last equality uses $\frac{\advr}{1-f}=\advr+\frac{\advr f}{1-f}$.
    This proves the first item. For the second, every honest validator has a
    chain of length $1$ at round $1$. From \Cref{lem.growth}, every honest
    validator adopts at round $v$ a chain of length at least $1+X([1,v-1])$.
    Since $v-1\ge\lambda$, \Cref{lem.typical} gives
    $X([1,v-1])>(1-\epsilon)f(v-1)$.
\end{proof}

Let $r\ge\lambda$. The first item bounds every chain of $\bt_r$ from above:
$\len{\ch}-1<\maxr fr$, whether or not any honest validator ever adopted $\ch$.
The second item bounds only the chains adopted by honest validators from below:
$\len{\ch}-1>(1-\epsilon)f(r-1)$. A chain that a client holds has no lower
bound, since the adversary may deliver a chain of any length. It does have an
upper bound. In particular, a height schedule with $\hgt(r)\ge\maxr fr$ for some
$r$ is not realizable at round $r$ in a typical execution.

\subsection{Counting lemmas}
\label{sec.counting}

The next three lemmas are the technical core. Throughout this subsection, let
$\ch\in\bt_r$ and let $\ch'$ be a chain adopted by an honest validator at round
$r$. Let $\bst$ be the last honest block of $\ch\meet\ch'$, let $\pos$ be its
height and let $\rst$ be the round at which it was computed, so
$\bst=\ch[\pos]=\ch'[\pos]$. Put $S=[\rst+1,r-1]$ and $S'=[\rst,r]$.
\Cref{lem.rounds} locates the rounds in which the blocks of $\ch[\pos+1{:}]$
and $\ch'[\pos+1{:}]$ were computed, \Cref{lem.honestlen} bounds
$\len{\ch'[\pos+1{:}]}$ from below, and \Cref{lem.count} bounds
$\len{\ch[\pos+1{:}]}$ from above.

\begin{lemma}\label{lem.rounds}
    $\rst\le r-1$. Every block of $\ch[\pos+1{:}]$ and every block of
    $\ch'[\pos+1{:}]$ was computed in a round of $S'$.
\end{lemma}
\begin{proof}
    From \Cref{def.tree} and \Cref{lem.membership}, every block of $\ch$ was
    computed in a round of $[0,r]$ and every block of $\ch'$ in a round of
    $[0,r-1]$. The block $\bst$ lies on $\ch'$ and was computed at round
    $\rst$, so $\rst\le r-1$. No predictions occur, so along a chain the round
    at which a block was computed is non-decreasing in the height. Hence every
    block of $\ch[\pos+1{:}]$ and every block of $\ch'[\pos+1{:}]$ was computed
    in a round of $[\rst,r]=S'$.
\end{proof}

\begin{lemma}\label{lem.honestlen}
    $\len{\ch'[\pos+1{:}]}\ge X(S)$.
\end{lemma}
\begin{proof}
    The block $\bst$ is honest, at height $\pos$, and was computed at round
    $\rst$. From
    \Cref{lem.fork}, some honest validator has a chain of length at least
    $\pos+1$ at round $\rst+1$. From \Cref{lem.rounds}, $\rst+1\le r$, so
    \Cref{lem.growth} applies with $v=\rst+1$ and $w=r$: every honest validator
    adopts at round $r$ a chain of length at least
    $\pos+1+X([\rst+1,r-1])=\pos+1+X(S)$. In particular
    $\len{\ch'[\pos+1{:}]}=\len{\ch'}-\pos-1\ge X(S)$.
\end{proof}

In words, every block of $\ch[\pos+1{:}]$ is adversarial, or is paired with a
distinct adversarial block of $\ch'$, or is one of at most $W(S)+1$ remaining
honest blocks.

\begin{lemma}\label{lem.count}
    $\len{\ch[\pos+1{:}]}\le Z(S')+W(S)+1$.
\end{lemma}
\begin{proof}
    Split the blocks of $\ch[\pos+1{:}]$ into three classes: $a$ adversarial
    blocks, $h_1$ honest blocks computed in a convergence opportunity of $S$,
    and $h_2$ remaining honest blocks. Then $\len{\ch[\pos+1{:}]}=a+h_1+h_2$.
    We bound $h_2$ first and $a+h_1$ second.

    From \Cref{lem.rounds}, every honest block of $\ch[\pos+1{:}]$ was
    computed in a round of $S'$. None was computed at round $\rst$: the block
    $\bst$ is honest and was computed at round $\rst$, and by
    \Cref{lem.honestone} a chain has at most one honest block per round. So
    every honest block of $\ch[\pos+1{:}]$ was computed in a round of
    $S\cup\{r\}$. The $h_2$ blocks are the honest blocks of $\ch[\pos+1{:}]$
    computed at round $r$ or in a successful round of $S$ that is not a
    convergence opportunity. From \Cref{lem.honestone} again, they were
    computed in distinct rounds. So at most one of them was computed at round
    $r$, and at most $W(S)$ of them were computed in successful rounds of $S$
    that are not convergence opportunities. Hence
    $h_2\le W(S)+1$.

    We claim that $a+h_1\le Z(S')$. Define a map from the $a+h_1$ blocks of the
    first two classes to adversarial blocks computed in rounds of $S'$. An adversarial $\ch[j]$ is mapped to itself. It was computed in a round of
    $S'$ by \Cref{lem.rounds}. An honest $\ch[j]$ computed in a convergence
    opportunity $u\in S$ is mapped to $\ch'[j]$. We check that $\ch'[j]$ is such
    a block. The honest validator that computed $\ch[j]$ sent a chain of length
    $j+1$ at round $u\le r-1$, so from \Cref{lem.fork} $\len{\ch'}>j$ and
    $\ch'[j]$ exists. Since $j>\pos$ and $\ch[j]$ is honest, $\ch'[j]\ne\ch[j]$ by the
    maximality of $\pos$, so $\ch'[j]$ is adversarial by \Cref{lem.unique}. It lies on $\ch'[\pos+1{:}]$, so
    from \Cref{lem.rounds} it was computed in a round of $S'$. The map is injective: two blocks at distinct heights of one chain are
    distinct, and $\ch[j]=\ch'[j']$ would force $j=j'$ since no copies occur,
    with $\ch[j]$ adversarial and $\ch[j']$ honest. Each adversarial block is
    the output of a distinct successful adversarial query, and $Z(S')$ counts
    the successful adversarial queries in the rounds of $S'$. Hence
    $a+h_1\le Z(S')$, and
    $\len{\ch[\pos+1{:}]}=a+h_1+h_2\le Z(S')+W(S)+1$.
\end{proof}

\subsection{Main statement}

\begin{theorem}\label{thm.prefix}
    In a typical execution, for every round $r$, every chain $\ch\in\bt_r$ and
    every chain $\ch'$ adopted by an honest validator at round $r$,
    \[
        \len{\ch\meet\ch'} > \len{\ch} - \max\left\{k_0,\
        \frac{\advrtwo\left(\len{\ch'}-\len{\ch}\right)+2f+1}{1-\advrtwo}\right\}.
    \]
\end{theorem}
\begin{proof}
    Let $\bst$, $\pos$, $\rst$, $S$ and $S'$ be as in \Cref{sec.counting}, and write $d=\len{\ch[\pos+1{:}]}=\len{\ch}-\pos-1$.
    Since $\ch\meet\ch'$ contains $\bst$, we have
    $\len{\ch\meet\ch'}\ge\pos+1=\len{\ch}-d$, so it suffices to bound $d$. If
    $d<k_0$ we are done, so assume $d\ge k_0=2(\lambda+2)f$.

    We first show that $|S|>\lambda$. The subsequence $\ch[\pos+1{:}]$ consists
    of $d$ consecutive blocks of $\ch$, so from \Cref{lem.span} they were
    computed in more than $d/(2f)\ge\lambda+2$ consecutive rounds. From
    \Cref{lem.rounds}, all these rounds lie in $S'=[\rst,r]$, which has $|S|+2$
    elements. Hence $|S|>\lambda$, and \Cref{lem.typical} applies to $S$ and to
    $S'$.

    From \Cref{lem.count,lem.typical},
    \[
        d < \left(\frac{\advr}{1-f}+\epsilon\right)f(|S|+2)
        + \left(\frac{f_2}{f}+\epsilon\right)f|S| + 1 .
    \]
    From \Cref{lem.rho}, $\frac{\advr}{1-f}+\epsilon<1$, so the first term is
    smaller than $\left(\frac{\advr}{1-f}+\epsilon\right)f|S|+2f$. Therefore
    \begin{equation}\label{eq.upper}
        d < \advrone f|S| + 2f + 1 .
    \end{equation}
    From \Cref{lem.honestlen,lem.typical},
    $\len{\ch'}-\pos-1\ge X(S)>(1-\epsilon)f|S|$, so
    $f|S|<(\len{\ch'}-\pos-1)/(1-\epsilon)$ and
    \[
        d < \advrtwo(\len{\ch'}-\pos-1)+2f+1 .
    \]
    Substitute $\len{\ch'}-\pos-1=(\len{\ch'}-\len{\ch})+d$. This gives
    $d<\advrtwo(\len{\ch'}-\len{\ch})+\advrtwo d+2f+1$, and $\advrtwo<1$ by
    \Cref{lem.rho}, so $d<\frac{\advrtwo(\len{\ch'}-\len{\ch})+2f+1}{1-\advrtwo}$.
\end{proof}

We call $\len{\ch'}-\len{\ch}$ the length deficit of $\ch$. \Cref{thm.prefix}
bounds the pruning depth by $\frac{\advrtwo}{1-\advrtwo}$ times the length
deficit plus the constant $\frac{2f+1}{1-\advrtwo}$. As $f$ and $\epsilon$ tend
to $0$ the coefficient of the length deficit tends to
\[
    \frac{\advr}{1-\advr} = \frac{t}{n-2t}.
\]
The bound is stated in terms
of $\len{\ch'}$, which a client does not know. The next corollary removes the
dependence on $\len{\ch'}$. The price is the width of the band of
\Cref{lem.length}.

\begin{corollary}\label{cor.round}
    In a typical execution, for every round $r$, every chain $\ch\in\bt_r$ and
    every chain $\ch'$ adopted by an honest validator at round $r$,
    \[
        \len{\ch\meet\ch'} > \len{\ch} - \max\left\{k_0,\
        \frac{\advrone\left[\maxr f(r+\lambda)-\len{\ch}+1\right]+\maxr(2f+1)}{1-f_2/f}\right\}.
    \]
\end{corollary}
\begin{proof}
    Let $\bst$, $\pos$, $\rst$ and $S$ be as in \Cref{sec.counting} and write $d=\len{\ch[\pos+1{:}]}=\len{\ch}-\pos-1$. As in the proof
    of \Cref{thm.prefix}, $\len{\ch\meet\ch'}\ge\pos+1=\len{\ch}-d$, so it suffices
    to bound $d$ by the maximum in the statement, and we may assume $d\ge k_0$,
    which gives $|S|>\lambda$ and \Cref{eq.upper}.

    The prefix $\ch[{:}\pos+1]$ is a chain of $\bt_{\rst}$, because its tip
    $\bst$ was computed at round $\rst$ and no predictions occur. From
    \Cref{lem.length}, $\pos<\maxr f\max(\rst,\lambda)\le\maxr f(\rst+\lambda)$,
    hence $\rst>\frac{\pos}{\maxr f}-\lambda$. Together with
    $|S|=r-\rst-1<r-\rst$ and \Cref{eq.upper} this gives
    \[
        d < \advrone f(r+\lambda) - \frac{\advrone}{\maxr}\pos + 2f+1 .
    \]
    Put $c=\advrone/\maxr$ and observe that $c<1$: $\advrone=(1-\epsilon)\advrtwo<1$ by
    \Cref{lem.rho}, and $\maxr=1+\advr+\slack>1$. Substitute $\pos=\len{\ch}-1-d$ and rearrange,
    \[
        d < \frac{\advrone f(r+\lambda)-c\,(\len{\ch}-1)+2f+1}{1-c} .
    \]
    Multiplying numerator and denominator by $\maxr$ gives
    \[
        d < \frac{\advrone\left[\maxr f(r+\lambda)-\len{\ch}+1\right]+\maxr(2f+1)}{\maxr-\advrone},
    \]
    and the definitions of $\maxr$ and $\advrone$ give
    $\maxr-\advrone=1-f_2/f$.
\end{proof}

Let $\hgt(r)$ be the height of $\ch$ and write $\pref(r)=\hgt(r)-\prn(r)$ for
the height of the prefix that remains after pruning $\prn(r)$ blocks. The
second term of the maximum in \Cref{cor.round} contains the additive terms
$\advrone\maxr f\lambda$ and $\maxr(2f+1)$, both of order $\lambda f$. Dropping
them and letting $f$ and $\epsilon$ tend to $0$, the bound of \Cref{cor.round}
becomes
\begin{gather*}
    \prn(r) = \frac{t}{n-t}\left(\frac{n}{n-t}fr-\hgt(r)\right),\\
    \pref(r) = \frac{n}{n-t}\left(\hgt(r)-\frac{t}{n-t}fr\right) .
\end{gather*}
Two regimes of this limiting form are worth naming. If $\hgt(r)\le\frac{t}{n-t}fr$, then
$\pref(r)\le0$: the chain is no longer than what the adversary alone can mine
in $r$ rounds, and nothing beyond the genesis block is certified. If $\hgt(r)\ge\frac{n}{n-t}fr-\frac{n-t}{t}k_0$, then $\prn(r)=k_0$. By
\Cref{lem.length} the chain is within $\frac{n-t}{t}k_0$ blocks of the longest
chain possible at round $r$. So at most that many blocks mined by round $r$ lie
off it, and the client prunes only the common prefix parameter $k_0$.

\subsection{Later rounds}

\Cref{cor.round} compares the client's chain with the honest chains of a
single round. A client that wakes up at round $r$ needs more: the prefix it
keeps must remain a prefix of every honest chain at every later round. The tool
is \Cref{thm.prefix} with the deficit at most $0$. Put
\[
    k_1 = \max\Big\{k_0,\ \frac{2f+1}{1-\advrtwo}\Big\}.
\]
If $\ch\in\bt_r$, $\ch'$ is adopted by an honest validator at round $r$ and
$\len{\ch}\ge\len{\ch'}$, then \Cref{thm.prefix} gives
\begin{equation}\label{eq.longer}
    \len{\ch\meet\ch'} > \len{\ch}-k_1 .
\end{equation}
In words, a chain at least as long as an honest chain forks from it by fewer
than $k_1$ blocks. The next lemma applies this twice: to the first honest chain
that drops the client's prefix, and then to the client's chain.

\begin{lemma}\label{lem.persist}
    Let $r$ be a round, let $\ch\in\bt_r$ and let $k\ge k_1$. If
    $\prune{\ch}{k}\preceq\ch'$ for every chain $\ch'$ adopted by an honest
    validator at round $r$, then $\prune{\ch}{k}\preceq\ch''$ for every chain
    $\ch''$ adopted by an honest validator at a round $r''>r$.
\end{lemma}
\begin{proof}
    Write $P=\prune{\ch}{k}$. If $k\ge\len{\ch}-1$ then $P$ is the genesis
    chain and there is nothing to prove, so assume $k<\len{\ch}-1$, which gives
    $\len{P}=\len{\ch}-k$. Suppose, towards a contradiction, that the claim
    fails. Let $s>r$ be the first
    round at which an honest validator $V$ adopts a chain $\ch_s$ with
    $P\not\preceq\ch_s$, and let $\ch_A$ be the chain $V$ adopted at round
    $s-1$. Then $P\preceq\ch_A$: by hypothesis if $s-1=r$, by the choice of $s$
    otherwise. At round $s$ the chain $V$ has is $\ch_A$ extended by at most
    one block. It contains $P$, so it differs from $\ch_s$, and since
    $\fn{maxvalid}$ returns a longest argument, $\len{\ch_s}\ge\len{\ch_A}$.

    First compare $\ch_s$ with $\ch_A$. Both $P$ and $\ch_s\meet\ch_A$ are
    prefixes of $\ch_A$ and $P\not\preceq\ch_s$, so
    $\len{\ch_s\meet\ch_A}<\len{P}$. On the other hand, $\ch_s\in\bt_{s-1}$ by
    \Cref{lem.membership} and $\len{\ch_s}\ge\len{\ch_A}$, so \Cref{eq.longer}
    at round $s-1$ gives $\len{\ch_s\meet\ch_A}>\len{\ch_s}-k_1\ge\len{\ch_s}-k$.
    Hence $\len{\ch_s}-k<\len{\ch}-k$, that is, $\len{\ch_s}<\len{\ch}$.

    Now compare $\ch$ with $\ch_s$. Since $\ch\in\bt_s$ and
    $\len{\ch}>\len{\ch_s}$, \Cref{eq.longer} at round $s$ gives
    $\len{\ch\meet\ch_s}>\len{\ch}-k_1\ge\len{\ch}-k=\len{P}$. Both $P$ and
    $\ch\meet\ch_s$ are prefixes of $\ch$, so
    $P\preceq\ch\meet\ch_s\preceq\ch_s$. This contradicts the choice of
    $\ch_s$.
\end{proof}

\begin{corollary}\label{cor.prune}
    Let $r$ be a round, let $\ch\in\bt_r$ and let $\prn(r)$ be the maximum
    appearing in \Cref{cor.round} for $\ch$ and $r$. If $\prn(r)\ge k_1$, then
    in a typical execution $\prune{\ch}{\prn(r)}\preceq\ch''$ for every chain
    $\ch''$ adopted by an honest validator at a round $r''\ge r$.
\end{corollary}
\begin{proof}
    For $r''=r$ this is \Cref{cor.round}:
    $\len{\ch\meet\ch''}>\len{\ch}-\prn(r)=\len{\prune{\ch}{\prn(r)}}$, and both
    are prefixes of $\ch$. For $r''>r$ it is \Cref{lem.persist} with
    $k=\prn(r)$.
\end{proof}

Since $\prn(r)\ge k_0$, the condition $\prn(r)\ge k_1$ holds whenever
$k_0\ge\frac{2f+1}{1-\advrtwo}$, that is, whenever
$\lambda\ge\frac{2f+1}{2f(1-\advrtwo)}-2$. We assume this from here on.
